\documentclass[10pt,a4paper]{amsart}
\usepackage[margin=19mm]{geometry}
\usepackage{amsmath,amssymb,mathtools}
\usepackage[scr=rsfs,cal=euler]{mathalfa}
\usepackage{xfrac}
\usepackage[unicode,hidelinks,bookmarks=false]{hyperref}
\newcommand{\ie}{i.e.\ }
\newcommand{\cf}{cf.\ }
\newcommand{\resp}{respectively\ }
\newcommand{\Subsection}{\subsection}
\providecommand{\nobreakdash}{\nobreak\hbox{-}}
\setlength{\emergencystretch}{2em}
\multlinegap0pt
\usepackage{mathtools}
\usepackage{amssymb}
\usepackage{bm}
\usepackage[shortlabels]{enumitem}
\usepackage{xcolor}
\newtheorem{lemma}{Lemma}
\newcommand{\EM}{\mathrm{EM}}
\newcommand{\e}{\mathrm{e}}
\newcommand{\eps}{\varepsilon}
\title{Long induced cycles in pseudorandom graphs}
\author{Sahar Diskin\textsuperscript{*}, Lyuben Lichev\textsuperscript{\textdagger}, Michael Krivelevich\textsuperscript{$\ddagger$}, Itay Markbreit\textsuperscript{\S}}
\date{Source: arXiv:2609.06176v1 (CC BY 4.0)}
\begin{document}
\maketitle

\noindent\emph{Source and licence.} Long induced cycles in pseudorandom graphs. Version arXiv:2609.06176v1 (CC BY 4.0), distributed by arXiv under the Creative Commons Attribution 4.0 International licence, http://creativecommons.org/licenses/by/4.0/, as recorded in the arXiv licence metadata for this exact version and shown on the abstract page https://arxiv.org/abs/2609.06176v1. Full licence notice: \texttt{licenses/arxiv-2609.06176-CC-BY-4.0.txt}.\par\smallskip

\noindent\textit{Editorial scope.} Counted unit: the article's lemma expansion (source0.tex lines 631--640). The article's complete original proof of it is delivered with it (source0.tex lines 642--731, one \texttt{proof} environment of 90 lines). No mathematics is written, repaired or re-derived: the statement and the proof are the article's own lines, in the article's own notation, and no retained line is substituted or re-flowed.  Retained uncounted as the source-local context the proof needs: lines 621--625.  Nothing was omitted from any retained span, so the package carries no omission record.  No retained line contains a citation, so the wrapper carries no bibliography and no bibliography entry.  No retained line contains a figure environment, an image-inclusion command or drawing code, so no image is delivered and the package declares no asset dependency.  Editorial formatting only, in addition: an AMS \texttt{amsart} preamble that restates the article's own declarations and macros used by the retained lines, namely the environments \texttt{lemma} on separate counters, together with the standard packages those lines need.  The retained environments print in this document as Lemma 1 in order of appearance, whereas the article prints the same items with the numbering of its own full document.  The complete native source of the article as distributed in the arXiv e-print for this exact version is bundled unchanged as \texttt{sources/209500/source0.tex}.
\begin{equation}\label{eq: theta-bounds}
    \delta_i,\theta_i\ll\beta\ll\eps,
    \qquad
    \theta_i\frac{n_{i-1}}{d_{i-1}} = \theta_i\frac{n}{d}\ge n^{4/5}.
\end{equation}

\begin{lemma}\label{l: expansion}
With probability $1-O(n^{-3})$, for every $i\in[k]$ and every subset $L\subseteq Y_i$ of size
\[
    \theta_i\frac{n_{i-1}}{d_{i-1}}\le |L|\le (1-\beta)\frac{n_{i-1}}{d_{i-1}},
\]
we have
\[
    |N_{G(i-1)}(L)|\ge (1-\beta)^2\left((1-\delta_i)d_{i-1}|L|-\frac{d_{i-1}^2|L|^2}{n_{i-1}}\right).
\]
\end{lemma}

\begin{proof}
Fix $i\in [k]$. 
We call a set $L\subseteq V(G(i-1))$ non-expanding if it violates the displayed inequality. We now bound from above the probability that $Y_i$ contains such a set. For any set $L$, we abbreviate $\ell\coloneqq |L|$.

Consider an ordered sequence $\tau=(v_1,\ldots,v_\ell)$ of distinct vertices. 
For all $j\in [\ell]$, write $L_j=\{v_1,\ldots,v_j\}$ and $N_j=N_{G(i-1)}(L_j)$. For $j\in \{2,\ldots, \ell\}$, given $v_1,\ldots,v_{j-1}$, call a vertex $v\in V(G(i-1))\setminus L_{j-1}$ \textit{bad} if it has fewer than
\begin{align}
 (1-\beta/4)\left(d_{i-1}-\frac{d_{i-1}(d_{i-1}+1)(j-1)}{n_{i-1}}\right)\label{eq: bad}   
\end{align}
neighbours in $V(G(i-1))\setminus(L_{j-1}\cup N_{j-1})$, and \textit{good} otherwise.

Suppose first that at most $\beta\ell/4$ vertices of $\tau$ are bad. Each good vertex increases the current open neighbourhood by at least the displayed quantity in \eqref{eq: bad}, except that vertices chosen later in the sequence may have been counted earlier as boundary vertices; the latter accounts for at most $\ell$ vertices. Since $\ell\le(1-\beta)n_{i-1}/d_{i-1}$, we get in this case
\begin{align}\label{eq: neigh bound}
    |N_{G(i-1)}(L_\ell)|\notag
    &\ge \left(\sum_{j=\beta\ell/4+1}^{\ell}(1-\beta/4)\left(d_{i-1}-\frac{d_{i-1}(d_{i-1}+1)(j-1)}{n_{i-1}}\right)\right)-\ell \\
    &\ge (1-\beta/4)^2d_{i-1}\ell-\ell-(1-\beta/4)\sum_{j=\beta\ell/4+1}^\ell\frac{d_{i-1}(d_{i-1}+1)(j-1)}{n_{i-1}}.
\end{align}
Now, observe that 
\[
    (1-\beta/4)\sum_{j=\beta\ell/4+1}^\ell\frac{d_{i-1}(d_{i-1}+1)(j-1)}{n_{i-1}}\le (1-\beta/4)\frac{d_{i-1}(d_{i-1}+1)}{n_{i-1}}\cdot \frac{\ell(\ell-1)}{2}\le  (1-\beta)^2\frac{d_{i-1}^2\ell^2}{n_{i-1}}.
\]
In addition, since $(1-\beta)^2\delta_id_{i-1}>1$, we have that 
\[
    (1-\beta/4)^2d_{i-1}\ell-\ell\ge (1-\beta)^2d_{i-1}\ell-(1-\beta)^2\delta_id_{i-1}\ell\ge (1-\beta)^2(1-\delta_{i})d_{i-1}\ell.
\]
Substituting these calculations back into \eqref{eq: neigh bound} yields
\begin{align*}
    |N_{G(i-1)}(L_\ell)|
    &\ge (1-\beta/4)^2d_{i-1}\ell-\ell-(1-\beta/4)\sum_{j=\beta\ell/4+1}^\ell\frac{d_{i-1}(d_{i-1}+1)(j-1)}{n_{i-1}}\\
    &\ge (1-\beta)^2\left((1-\delta_{i})d_{i-1}\ell-\frac{d_{i-1}^2\ell^2}{n_{i-1}}\right).
\end{align*}
Thus, every ordering of a non-expanding set has at least $\beta\ell/4$ bad positions.

We next bound the number of choices for a bad vertex at a fixed position $j$. Let $X=L_{j-1}\cup N_{j-1}$, let $U=V(G(i-1))\setminus X$, and let $B$ be the set of bad vertices. Since $G(i-1)$ is $(\delta_{i-1},d_{i-1})$-nearly regular and $|V(G(i-1))|=(1\pm2\delta_{i-1})n_{i-1}$, we have
\begin{align}
    |U|\ge (1-2\delta_{i-1})n_{i-1}-((1+\delta_{i-1})d_{i-1}+1)(j-1)\ge \left(\beta/2-3\delta_{i-1}\right)n_{i-1}\ge \beta n_{i-1}/3,\label{eq: u and beta}
\end{align}
where the second inequality uses that $j\le\ell\le(1-\beta)n_{i-1}/d_{i-1}$, and the last inequality uses that $\delta_{i-1}\ll \beta$. 
Further, note that
\begin{align}
    d_{i-1}-\frac{d_{i-1}(d_{i-1}+1)(j-1)}{n_{i-1}}\le \frac{d_{i-1}|U|}{n_{i-1}}+3\delta_{i-1}d_{i-1}\le (1+\beta/20)\frac{d_{i-1}|U|}{n_{i-1}}\label{eq: bounding inner term}
\end{align}
where the last inequality uses \eqref{eq: u and beta} and that $\delta_{i-1}\ll \beta$. 
Hence, by definition of the bad vertices,
\begin{align}
    e(B,U)&\le (1-\beta/4)\left(d_{i-1}-\frac{d_{i-1}(d_{i-1}+1)(j-1)}{n_{i-1}}\right)|B|\notag\\
    &\le (1-\beta/4)(1+\beta/20)\frac{d_{i-1}|U|}{n_{i-1}}|B|\le (1-\beta/10)\frac{d_{i-1}|U|}{n_{i-1}}|B|,\label{eq: upper bound on e(b,u)}
\end{align}
where the second inequality follows from \eqref{eq: bounding inner term}.
On the other hand, by the property $\EM(n,(1\pm \gamma)d,\lambda)$ of $G(i-1)$,
\[
e(B,U)\ge \frac{(1-\gamma)^2d|B||U|}{(1+\gamma)n}-\frac{1+\gamma}{1-\gamma}\lambda\sqrt{|B||U|},
\]
and thus by \eqref{eq: upper bound on e(b,u)},
\begin{align}\label{eq: lambda bound B U}
    \lambda\sqrt{|B||U|}\ge \frac{(1-\gamma)^3d_{i-1}|B||U|}{(1+\gamma)^2n_{i-1}}-\frac{1-\gamma}{1+\gamma}\cdot \left(1-\frac{\beta}{10}\right)\frac{d_{i-1}|U|}{n_{i-1}}|B|\ge \frac{\beta d_{i-1}}{20n_{i-1}}|B||U|,
\end{align}
where the last inequality holds since $\gamma\ll \beta$. By combining \eqref{eq: u and beta} and \eqref{eq: lambda bound B U}, we obtain
\begin{equation}\label{eq: bad-choices}
    |B|\le \left(\frac{20\lambda}{\beta d_{i-1}}\right)^2\frac{n_{i-1}^2}{|U|}\le \left(\frac{20}{\beta}\right)^3\left(\frac{\lambda}{d_{i-1}}\right)^2n_{i-1}.
\end{equation}

By using that $\tbinom{\ell}{\beta\ell/4}\le (4\e/\beta)^{\beta\ell/4}$, it follows that the number of ordered non-expanding $\ell$-tuples is at most
\[
    \binom{\ell}{\beta\ell/4}
    \bigg(\bigg(\frac{20}{\beta}\bigg)^3\left(\frac{\lambda}{d_{i-1}}\right)^2n_{i-1}\bigg)^{\beta\ell/4}
    n_{i-1}^{(1-\beta/4)\ell}
    \le
    \bigg(\bigg(\bigg(\frac{20}{\beta}\bigg)^4\left(\frac{\lambda}{d_{i-1}}\right)^2\bigg)^{\beta/4}n_{i-1}\bigg)^\ell.
\]
Multiplying by $p_i^{\ell}$
%$\left(\frac{1+\eps}{d_{i-1}}\right)^\ell$ 
and dividing by $\ell!$, the probability that $Y_i$ contains a non-expanding $\ell$-set is at most
\[
    \left[
    \left(\bigg(\frac{20}{\beta}\bigg)^4\left(\frac{\lambda}{d_{i-1}}\right)^2\right)^{\beta/4}
    \frac{(1+\eps)\e n_{i-1}}{d_{i-1}\ell}
    \right]^\ell.
\]
Since $\ell\ge\theta_i n_{i-1}/d_{i-1}$ and $\theta_i=(\lambda/d_{i-1})^{\beta/8}\ll \beta\ll 1$, the expression in the square brackets above is at most
\[
    (1+\eps)\e\cdot \left(\frac{20}{\beta}\right)^\beta \left(\frac{\lambda}{d_{i-1}}\right)^{3\beta/8}\le (1+\eps)\e\cdot \frac{20}{\beta} \left(\frac{\lambda}{d_{i-1}}\right)^{3\beta/8}\le 100 \frac{\theta_i^3}{\beta}\le \frac{1}{2}.
\]
%which is at most $1/2$ since $\theta_i\ll \beta\ll 1$. 
Summing over all admissible $\ell$ gives a failure probability at most
\[
    \sum_{\ell\ge \theta_i n_{i-1}/d_{i-1}}2^{-\ell}=O(n^{-4})
\]
by \eqref{eq: theta-bounds}, and the lemma follows by a union bound over $i\in [k]$.
\end{proof}

\end{document}
